\chapter{Selbsttests und Übungen}
\label{cha:Tests}

Dieses Kapitel enthält Selbsttests zur Überprüfung Ihrer Deutschkenntnisse.

\section{A1-A2 Grundlagen Test}

\bigskip
\noindent\textbf{Aufgabe 1: Konjugieren Sie die Verben im Präsens}

\begin{enumerate}
    \item Ich \_\_\_\_\_\_\_ (arbeiten) bei Siemens.
    \item Er \_\_\_\_\_\_\_ (wohnen) in Berlin.
    \item Wir \_\_\_\_\_\_\_ (fahren) nach Hamburg.
    \item Sie \_\_\_\_\_\_\_ (sprechen) Englisch.
    \item Ich \_\_\_\_\_\_\_ (haben) zwei Kinder.
\end{enumerate}

\bigskip
\textbf{Lösung 1:} arbeite, wohnt, fahren, sprechen, habe

\bigskip
\noindent\textbf{Aufgabe 2: Bestimmen Sie den Artikel (der/die/das)}

\begin{enumerate}
    \item \_\_\_\_\_\_\_ Buch
    \item \_\_\_\_\_\_\_ Schule
    \item \_\_\_\_\_\_\_ Haus
    \item \_\_\_\_\_\_\_ Frau
    \item \_\_\_\_\_\_\_ Kind
\end{enumerate}

\bigskip
\textbf{Lösung 2:} das, die, das, die, das

\bigskip
\noindent\textbf{Aufgabe 3: Bilden Sie den Plural}

\begin{enumerate}
    \item das Kind → die \_\_\_\_\_\_\_
    \item der Mann → die \_\_\_\_\_\_\_
    \item das Haus → die \_\_\_\_\_\_\_
    \item die Frau → die \_\_\_\_\_\_\_
    \item der Student → die \_\_\_\_\_\_\_
\end{enumerate}

\bigskip
\textbf{Lösung 3:} Kinder, Männer, Häuser, Frauen, Studenten

\bigskip
\noindent\textbf{Aufgabe 4: Perfekt bilden}

\begin{enumerate}
    \item Ich \_\_\_\_\_\_\_ gestern \_\_\_\_\_\_\_ (kommen).
    \item Er \_\_\_\_\_\_\_ das Buch \_\_\_\_\_\_\_ (lesen).
    \item Wir \_\_\_\_\_\_\_ nach Berlin \_\_\_\_\_\_\_ (fahren).
    \item Sie \_\_\_\_\_\_\_ viel \_\_\_\_\_\_\_ (arbeiten).
\end{enumerate}

\bigskip
\textbf{Lösung 4:} bin gekommen, hat gelesen, sind gefahren, haben gearbeitet

\section{B1 Mittelstufe Test}

\bigskip
\noindent\textbf{Aufgabe 1: Konjunktiv II}

\begin{enumerate}
    \item Wenn ich Zeit \_\_\_\_\_\_\_, \_\_\_\_\_\_\_ ich dich besuchen.
    \item Er \_\_\_\_\_\_\_ gern reisen.
    \item \_\_\_\_\_\_\_ du mir bitte helfen?
    \item Ich \_\_\_\_\_\_\_ gern in Berlin wohnen.
\end{enumerate}

\bigskip
\textbf{Lösung 1:} hätte ... würde, würde, Könntest, würde

\bigskip
\noindent\textbf{Aufgabe 2: Relativsätze}

\begin{enumerate}
    \item Das ist der Mann, \_\_\_\_\_\_\_ ich gestern getroffen habe.
    \item Die Schule, \_\_\_\_\_\_\_ ich besuche, ist sehr groß.
    \item Ich kenne jemanden, \_\_\_\_\_\_\_ Deutsch spricht.
    \item Das Buch, \_\_\_\_\_\_\_ ich gelesen habe, war interessant.
\end{enumerate}

\bigskip
\textbf{Lösung 2:} den, die, der, das

\bigskip
\noindent\textbf{Aufgabe 3: Passiv}

\begin{enumerate}
    \item Das Haus \_\_\_\_\_\_\_ (bauen) gerade.
    \item Die Aufgabe \_\_\_\_\_\_\_ (erledigen) schon.
    \item Er \_\_\_\_\_\_\_ (operieren) gestern.
    \item Die Arbeit \_\_\_\_\_\_\_ (verteilen) auf alle.
\end{enumerate}

\bigskip
\textbf{Lösung 3:} wird gebaut, ist erledigt, wurde operiert, wird verteilt

\bigskip
\noindent\textbf{Aufgabe 4: Präpositionen}

\begin{enumerate}
    \item Ich warte \_\_\_\_\_\_\_ dich \_\_\_\_\_\_\_ dem Bahnhof.
    \item Er kommt \_\_\_\_\_\_\_ dem Meeting \_\_\_\_\_\_\_ 10 Uhr.
    \item Sie ist \_\_\_\_\_\_\_ ihrer Arbeit sehr zufrieden.
    \item Wir sprechen \_\_\_\_\_\_\_ das Problem.
\end{enumerate}

\bigskip
\textbf{Lösung 4:} vor, um, mit, über

\section{B2-C1 Fortgeschrittenen Test}

\bigskip
\noindent\textbf{Aufgabe 1: Konjunktiv I (indirekte Rede)}

\begin{enumerate}
    \item Er sagte, er \_\_\_\_\_\_\_ krank.
    \item Sie meinte, sie \_\_\_\_\_\_\_ das nicht machen.
    \item Der Journalist schrieb, die Situation \_\_\_\_\_\_\_ sich verbessern.
    \item Man munkelt, er \_\_\_\_\_\_\_ ein Genie sein.
\end{enumerate}

\bigskip
\textbf{Lösung 1:} sei, wolle, habe sich verbessert, sei

\bigskip
\noindent\textbf{Aufgabe 2: Partizipialattribute}

\begin{enumerate}
    \item das \_\_\_\_\_\_\_ Buch (lesen)
    \item die \_\_\_\_\_\_\_ Frau (arbeiten)
    \item das \_\_\_\_\_\_\_ Problem (lösen)
    \item die \_\_\_\_\_\_\_ Antwort (geben)
\end{enumerate}

\bigskip
\textbf{Lösung 2:} gelesene, arbeitende, gelöste, gegebene

\bigskip
\noindent\textbf{Aufgabe 3: Funktionsverbgefüge}

\begin{enumerate}
    \item eine Entscheidung \_\_\_\_\_\_\_ (treffen)
    \item in \_\_\_\_\_\_\_ kommen (Bewegung)
    \item Zur \_\_\_\_\_\_\_ kommen (Entscheidung)
    \item unter \_\_\_\_\_\_\_ stellen (Frage)
    \item einen \_\_\_\_\_\_\_ machen (Vorschlag)
\end{enumerate}

\bigskip
\textbf{Lösung 3:} treffen, Gang, Entscheidung, Frage, Vorschlag

\bigskip
\noindent\textbf{Aufgabe 4: Nebensätze mit Subjunktoren}

\begin{enumerate}
    \item \_\_\_\_\_\_\_ ich Zeit habe, komme ich.
    \item \_\_\_\_\_\_\_ er kam, begann die Feier.
    \item \_\_\_\_\_\_\_ ich auch versuchte, es gelang mir nicht.
    \item \_\_\_\_\_\_\_ sie arbeitet, verdient sie wenig.
\end{enumerate}

\bigskip
\textbf{Lösung 4:} Wenn, Als, Obwohl, Während

\chapterend

\chapter{Vokabelkarten (Flashcards)}

Dieses Kapitel enthält wichtige Vokabelgruppen zum Lernen.

\section{Verben mit Akkusativ}

\bigskip
\begin{center}
\begin{tabular}{|l|p{6cm}|}
\hline
Verb & Bedeutung \\ \hline
sehen & to see \\ \hline
hören & to hear \\ \hline
verstehen & to understand \\ \hline
kennen & to know (person) \\ \hline
brauchen & to need \\ \hline
suchen & to search for \\ \hline
finden & to find \\ \hline
haben & to have \\ \hline
lieben & to love \\ \hline
hassen & to hate \\ \hline
rufen & to call \\ \hline
legen & to lay (down) \\ \hline
stellen & to put (standing) \\ \hline
setzen & to put (seated) \\ \hline
\end{tabular}
\end{center}

\section{Verben mit Dativ}

\bigskip
\begin{center}
\begin{tabular}{|l|p{6cm}|}
\hline
Verb & Bedeutung \\ \hline
helfen & to help \\ \hline
danken & to thank \\ \hline
gehören & to belong to \\ \hline
trauen & to trust \\ \hline
antworten & to answer \\ \hline
glauben & to believe \\ \hline
zustimmen & to agree \\ \hline
widersprechen & to contradict \\ \hline
begegnen & to meet (by chance) \\ \hline
fehlen & to miss / to be missing \\ \hline
gefallen & to like (lit. to be pleasing) \\ \hline
schmecken & to taste \\ \hline
gratulieren & to congratulate \\ \hline
\end{tabular}
\end{center}

\section{Verben mit Dativ und Akkusativ}

\bigskip
\begin{center}
\begin{tabular}{|l|l|l|}
\hline
Verb & Dativ & Akkusativ \\ \hline
geben & mir & das Buch \\ \hline
zeigen & mir & das Bild \\ \hline
sagen & mir & die Wahrheit \\ \hline
erzählen & mir & eine Geschichte \\ \hline
kaufen & mir & ein Auto \\ \hline
bringen & mir & die Post \\ \hline
schicken & mir & einen Brief \\ \hline
empfehlen & mir & ein Restaurant \\ \hline
erklären & mir & die Grammatik \\ \hline
versprechen & mir & die Wahrheit \\ \hline
\end{tabular}
\end{center}

\section{Falsche Freunde - Quick Reference}

\bigskip
\begin{center}
\begin{tabular}{|l|l|l|}
\hline
Deutsch & Falsch & Richtig \\ \hline
aktuell & actual & current \\ \hline
Gift & gift & poison \\ \hline
Billion & billion & trillion \\ \hline
fast & fast & almost \\ \hline
eventuell & eventful & possibly \\ \hline
konsequent & consequent & consistent \\ \hline
sensitiv & sensitiv & sensitive \\ \hline
sympathisch & sympathisch & likeable \\ \hline
praktisch & practical & practical, handy \\ \hline
real & real & realistic, real \\ \hline
\end{tabular}
\end{center}

\section{Checklisten für Grammatik}

\bigskip
\noindent\textbf{Vor dem Test - Checkliste:}

\begin{itemize}
    \item [ ] Alle Personalpronomen (Nominativ, Akkusativ, Dativ)
    \item [ ] Zeiten: Präsens, Präteritum, Perfekt, Plusquamperfekt, Futur I/II
    \item [ ] Artikel: bestimmt, unbestimmt, Negativartikel
    \item [ ] Pluralbildung der Nomen
    \item [ ] Deklination: Nominativ, Akkusativ, Dativ, Genitiv
    \item [ ] Adjektivdeklination
    \item [ ] Konjunktionen und Subjunktoren
    \item [ ] Modalverben
    \item [ ] Passiv
    \item [ ] Relativsätze
    \item [ ] Präpositionen mit Akkusativ/Dativ/Genitiv
    \item [ ] Konjunktiv I und II
\end{itemize}

\bigskip
\noindent\textbf{Wochenplan für Grammatik-Lernen:}

\bigskip
\begin{center}
\begin{tabular}{|l|l|}
\hline
Tag & Thema \\ \hline
Montag & Verben konjugieren (Präsens) \\ \hline
Dienstag & Artikel und Plural \\ \hline
Mittwoch & Satzbau und Wortstellung \\ \hline
Donnerstag & Zeiten (Perfekt/Präteritum) \\ \hline
Freitag & Präpositionen \\ \hline
Samstag & Wiederholung und Übungen \\ \hline
Sonntag & Aktiver Wortschatz \\ \hline
\end{tabular}
\end{center}

\chapterend